% Figure: blackbody spectral radiance (Planck's law) at three temperatures,
% showing the peak shift to shorter wavelength with rising temperature
% (Wien's law) and the absence of the ultraviolet catastrophe that the
% classical Rayleigh-Jeans law predicted. Plotted as relative spectral
% radiance versus wavelength.
\begin{figure}[htbp]
\centering
\begin{tikzpicture}
\begin{axis}[
  width=12cm, height=7.5cm,
  xlabel={wavelength $\lambda$ ($\mu$m)},
  ylabel={spectral radiance (relative)},
  xmin=0, xmax=3.0, ymin=0, ymax=1.05,
  axis lines=left,
  legend pos=north east,
  legend cell align=left,
  tick label style={font=\scriptsize},
  label style={font=\small},
  legend style={font=\scriptsize},
  clip=false,
]
% Planck law shape, normalised to unit peak for each temperature.
% B(lambda) ~ 1/lambda^5 * 1/(exp(c2/(lambda T)) - 1). Use c2 = 14388 um.K.
% Normalise each curve by its own peak so all peaks reach ~1.
% 5000 K: peak at lambda = 2898/5000 = 0.58 um.
\addplot[engred, very thick, domain=0.30:3.0, samples=120]
  {(1/x^5)/(exp(14388/(x*5000))-1) / ((1/0.58^5)/(exp(14388/(0.58*5000))-1))};
\addlegendentry{$T = 5000$ K}
% 4000 K: peak at 0.72 um.
\addplot[engorange, very thick, domain=0.30:3.0, samples=120]
  {(1/x^5)/(exp(14388/(x*4000))-1) / ((1/0.72^5)/(exp(14388/(0.72*4000))-1))};
\addlegendentry{$T = 4000$ K}
% 3000 K: peak at 0.97 um.
\addplot[engblue, very thick, domain=0.30:3.0, samples=120]
  {(1/x^5)/(exp(14388/(x*3000))-1) / ((1/0.966^5)/(exp(14388/(0.966*3000))-1))};
\addlegendentry{$T = 3000$ K}
% visible band marker
\draw[enggray, dashed] (axis cs:0.40,0) -- (axis cs:0.40,1.0);
\draw[enggray, dashed] (axis cs:0.70,0) -- (axis cs:0.70,1.0);
\node[enggray, font=\scriptsize, anchor=south] at (axis cs:0.55,1.0) {visible};
\end{axis}
\end{tikzpicture}
\caption[Blackbody spectral radiance from Planck's law]{Spectral radiance of
a blackbody at three temperatures, each curve normalised to its own peak.
The peak wavelength shifts toward shorter wavelength as temperature rises
(Wien's displacement law, $\lambda_{\max} T = 2898\,\mu\text{m}\cdot\text{K}$),
which is why a heated filament glows first dull red, then orange, then white.
The curves fall to zero at short wavelength: Planck's quantisation of the
field removes the ultraviolet divergence that the classical equipartition
result predicted. The relative areas under the curves grow as $T^4$
(Stefan-Boltzmann).}
\label{fig:vol04ch12:blackbody-curves}
\end{figure}
